\input{template.tex}

\ifPDFTeX
  \renewcommand{\rmdefault}{cmr}
\fi

\def\SheetNo{6}
\setlength{\emergencystretch}{2em}

\begin{document}
\sheettitle{Разбиения + комбинаторика.}

\begin{center}\textbf{Решения}\end{center}

\noindent В задачах о разбиениях числа на натуральные слагаемые порядок слагаемых не учитывается. Разбиение нуля - пустая сумма; разбиений отрицательного числа нет.

\begin{problems}

\item На доске написано несколько натуральных чисел $a_0,a_1,\ldots,a_n$. На второй доске записывают числа: $b_0$ - сколько всего чисел на первой доске, $b_1$ - сколько там чисел, больших единицы, $b_2$ - сколько чисел, больших двойки, и так далее, пока получаются положительные числа. Нули не записывают. На третьей доске записывают числа $c_0,c_1,\ldots$, построенные по числам второй доски по тому же правилу. Докажите, что наборы чисел на первой и третьей досках совпадают с учётом кратностей; порядок чисел не важен.
\par\smallskip\noindent\textbf{Решение.} Упорядочим числа $a_0,\ldots,a_n$ по невозрастанию и изобразим диаграмму Юнга: каждому числу соответствует столбец такой высоты. Число $b_j$ равно количеству столбцов высоты хотя бы $j+1$, то есть длине $(j+1)$-й строки диаграммы.
\par Следовательно, переход от набора $a$ к набору $b$ означает транспонирование диаграммы Юнга. Переход от $b$ к $c$ транспонирует её ещё раз. Два транспонирования возвращают исходную диаграмму, поэтому наборы $a$ и $c$ совпадают с учётом кратностей.\par

\item В классе учатся $32$ человека. Для участия в стритболе они организовали $33$ различные команды по три человека в каждой. Докажите, что найдутся две команды, у которых ровно один общий участник.
\par\smallskip\noindent\textbf{Решение.} Предположим, что требуемой пары команд нет. Тогда две различные команды либо не пересекаются, либо имеют ровно двух общих участников.
\par Разобьём команды на группы. Выберем команду $A$ и объединим в одну группу её саму и все команды, пересекающиеся с ней. Любая другая команда $T$ не пересекается с $A$. При этом каждая команда выбранной группы имеет не больше одного участника вне $A$, поэтому её пересечение с $T$ содержит не больше одного человека. По предположению пересечение из одного человека запрещено, значит, $T$ не пересекается ни с одной командой группы. Удалим всю группу и повторим построение. Таким образом, наборы участников разных групп попарно не пересекаются.
\par Покажем, что внутри каждой группы команд не больше, чем участников. Если команда всего одна, это очевидно. Иначе возьмём
\[
A=\{a,b,c\},\qquad B=\{a,b,d\}.
\]
Любые две команды группы пересекаются: каждая содержит хотя бы двух участников из $A$, поэтому они имеют хотя бы одного общего участника, а значит, ровно двух.
\par Если все команды группы содержат пару $\{a,b\}$, то каждой соответствует свой третий участник. При $t$ командах участников $t+2$, и нужная оценка выполнена.
\par Иначе существует команда $C$, не содержащая эту пару. Для пересечения по двум участникам и с $A$, и с $B$ она должна иметь вид $\{a,c,d\}$ или $\{b,c,d\}$. Без ограничения общности $C=\{a,c,d\}$. У любой другой команды не может быть участника вне $\{a,b,c,d\}$: тогда два оставшихся должны были бы одновременно входить в $A$, $B$ и $C$, но $A\cap B\cap C=\{a\}$. Следовательно, все команды группы лежат в множестве из четырёх участников, и их не больше $\binom43=4$.
\par Итак, в каждой группе команд не больше, чем участников. Группы используют непересекающиеся наборы участников, поэтому всего команд не больше $32$. Это противоречит наличию $33$ различных команд.\par

\item Вершины правильного стоугольника покрашены в $10$ цветов. Докажите, что можно выбрать четыре его вершины, являющиеся вершинами прямоугольника и покрашенные не более чем в два различных цвета.
\par\smallskip\noindent\textbf{Решение.} Пронумеруем вершины $a_1,\ldots,a_{100}$. Пары
\[
(a_i,a_{i+50}),\qquad i=1,\ldots,50,
\]
задают $50$ диаметров описанной окружности. Концы любых двух различных диаметров образуют прямоугольник.
\par Если есть два диаметра с одноцветными концами, их четыре конца используют не больше двух цветов, и задача решена.
\par Иначе одноцветный диаметр не более чем один, поэтому хотя бы $49$ диаметров имеют разноцветные концы. Неупорядоченных пар различных цветов всего $\binom{10}{2}=45$. По принципу Дирихле у двух диаметров пары цветов совпадут. Их концы дадут искомый прямоугольник.\par

\item Тест по английскому языку сдавали $10$ школьников. Известно, что любые пять школьников вместе ответили на все вопросы, а любые четыре вместе ответили не на все вопросы. При каком наименьшем количестве вопросов теста такое могло случиться? Группа ответила на вопрос, если на него ответил хотя бы один участник группы.
\par\smallskip\noindent\textbf{Решение.} Для каждой четвёрки школьников существует вопрос, на который никто из неё не ответил. При этом на один вопрос не могут не ответить две различные четвёрки: в их объединении хотя бы пять школьников, и тогда найдётся пятёрка, не ответившая на этот вопрос, вопреки условию. Поэтому разным четвёркам необходимы разные вопросы. Получаем нижнюю оценку
\[
\binom{10}{4}=210.
\]
\par Построим пример. Сопоставим каждой четвёрке отдельный вопрос. На него не отвечают ровно эти четверо, а остальные шестеро отвечают. Каждая четвёрка не отвечает на свой вопрос. Любые пять школьников отвечают вместе на каждый вопрос, поскольку не могут все входить в четвёрку не ответивших на него. Следовательно, условия выполнены, и ответ равен
\[
210.
\]\par

\item Пусть $0\le n\le k\ell$. Докажите, что количество разбиений числа $n$ в сумму не более чем $k$ слагаемых, каждое из которых не превосходит $\ell$, равно количеству разбиений числа $k\ell-n$ в сумму не более чем $k$ слагаемых, каждое из которых не превосходит $\ell$.
\par\smallskip\noindent\textbf{Решение.} Изобразим разбиение диаграммой Юнга, помещённой в прямоугольник из $k$ строк и $\ell$ столбцов. Её площадь равна $n$. Возьмём дополнение диаграммы до всего прямоугольника и повернём его на $180^\circ$. Получится диаграмма Юнга площади $k\ell-n$, также помещающаяся в этот прямоугольник.
\par Формально, дополнив разбиение нулями, запишем его как
\[
\ell\ge\lambda_1\ge\cdots\ge\lambda_k\ge0.
\]
Ему соответствует разбиение с частями
\[
\mu_i=\ell-\lambda_{k+1-i},\qquad i=1,\ldots,k,
\]
после удаления нулевых частей. Сумма частей равна $k\ell-n$. Повторное преобразование возвращает исходное разбиение, поэтому это биекция.\par

\item Петя написал количество разбиений числа $n$ в сумму не более чем $k$ слагаемых, каждое из которых не превосходит $\ell$. Вася написал количество разбиений числа $n$ в сумму не более чем $k$ слагаемых, каждое из которых не превосходит $\ell-1$. Маша написала количество разбиений числа $n-\ell$ в сумму не более чем $k-1$ слагаемых, каждое из которых не превосходит $\ell$. Докажите, что сумма Машиного и Васиного чисел равна числу Пети.
\par\smallskip\noindent\textbf{Решение.} Разделим все разбиения, посчитанные Петей, на два непересекающихся класса. Если слагаемого $\ell$ нет, все слагаемые не превосходят $\ell-1$; это в точности разбиения, посчитанные Васей.
\par Если слагаемое $\ell$ есть, удалим один его экземпляр. Останется разбиение числа $n-\ell$ на не более чем $k-1$ слагаемых, каждое не больше $\ell$. Это разбиения Маши. Обратное преобразование --- добавить одно слагаемое $\ell$ --- однозначно. Следовательно, число Пети равно сумме чисел Васи и Маши.\par

\item Три пирата Джо, Билл и Том нашли клад, содержащий $80$ одинаковых золотых монет, и хотят разделить их так, чтобы каждому досталось не менее $15$ монет. Сколько существует способов это сделать?
\par\smallskip\noindent\textbf{Решение.} Пусть $x_1,x_2,x_3$ --- количества монет у Джо, Билла и Тома соответственно. Требуется посчитать решения
\[
x_1+x_2+x_3=80,\qquad x_i\ge15.
\]
Положим $y_i=x_i-15$. Тогда
\[
y_1+y_2+y_3=35,\qquad y_i\ge0.
\]
По формуле шаров и перегородок число решений равно
\[
\binom{37}{2}=666.
\]
Пираты различимы, поэтому дополнительно делить на факториалы не требуется.\par

\item Петя подсчитал количество всех возможных $m$-буквенных слов из букв $\mathrm{T}$, $\mathrm{O}$, $\mathrm{W}$ и $\mathrm{N}$, в каждом из которых букв $\mathrm{T}$ и $\mathrm{O}$ поровну. Вася подсчитал количество всех возможных $2m$-буквенных слов из букв $\mathrm{T}$ и $\mathrm{O}$, в каждом из которых этих букв поровну. У кого слов получилось больше? Слово - любая последовательность букв.
\par\smallskip\noindent\textbf{Решение.} Построим биекцию между словами, посчитанными Петей и Васей. Каждую букву Петиного слова заменим парой букв:
\[
\mathrm{T}\mapsto\mathrm{TT},\qquad
\mathrm{O}\mapsto\mathrm{OO},\qquad
\mathrm{W}\mapsto\mathrm{TO},\qquad
\mathrm{N}\mapsto\mathrm{OT}.
\]
Получится слово длины $2m$. Разность количеств букв $\mathrm{T}$ и $\mathrm{O}$ при замене удваивается: первая замена даёт вклад $2$, вторая --- $-2$, а последние две --- $0$. Поэтому в новом слове этих букв поровну, то есть по $m$.
\par Обратно, разобьём Васино слово на последовательные пары букв. Каждую пару однозначно заменим одной буквой по той же таблице. Разность количеств букв $\mathrm{T}$ и $\mathrm{O}$ в восстановленном слове равна половине исходной разности, то есть нулю. Следовательно, получится Петино слово. Преобразования взаимно обратны, поэтому слов у Пети и Васи поровну, а их количество равно
\[
\binom{2m}{m}.
\]\par

\item Найдите количество восьмизначных чисел, произведение цифр которых равно $1400$. Ответ представьте в виде целого числа.
\par\smallskip\noindent\textbf{Решение.} Разложим произведение:
\[
1400=2^3\cdot5^2\cdot7.
\]
Нуля среди цифр быть не может. Цифры $3,6,9$ также невозможны, поскольку произведение не делится на $3$. Два множителя $5$ дают ровно две цифры $5$, а множитель $7$ --- одну цифру $7$. Оставшийся множитель $2^3$ можно получить тремя способами: цифрой $8$, цифрами $4,2$ или цифрами $2,2,2$. Остальные позиции заполняются единицами.
\par Таким образом, возможны ровно три набора цифр:
\[
\begin{array}{c|c}
\text{Набор цифр}&\text{Число перестановок}\\ \hline
1,1,1,1,5,5,7,8 & \dfrac{8!}{4!\,2!}=840\\[0.8em]
1,1,1,2,4,5,5,7 & \dfrac{8!}{3!\,2!}=3360\\[0.8em]
1,1,2,2,2,5,5,7 & \dfrac{8!}{2!\,3!\,2!}=1680
\end{array}
\]
Все цифры ненулевые, поэтому каждая перестановка задаёт восьмизначное число. Ответ:
\[
840+3360+1680=5880.
\]\par

\item Паша, Оля и Саша играли в тетрис. Победитель очередного раунда получал одно очко, остальные - по нулю. Игра заканчивается, как только кто-то наберёт $101$ очко. У Оли всё ещё было $0$ очков, когда один из двух остальных участников впервые набрал $100$. Но затем Оля собралась с силами и победила, причём в конце игры у Паши и Саши было по $100$ очков. Сколько существует последовательностей побед в раундах, при которых такое могло произойти?
\par\smallskip\noindent\textbf{Решение.} Обозначим победителей раундов буквами П, С, О. В конце у Паши и Саши по $100$ очков, у Оли $101$, причём последний раунд обязательно выиграла Оля. Удалим эту последнюю букву О. Останется последовательность длины $300$ со $100$ буквами каждого вида.
\par Условие означает, что до первой победы Оли либо Паша, либо Саша уже набрал $100$ очков. Рассмотрим первый случай: все $100$ букв П стоят раньше всех букв О. Если удалить буквы С, останется единственная последовательность $\text{П}^{100}\text{О}^{100}$. Для её дополнения выбираем $100$ позиций букв С среди $300$; получаем $\binom{300}{100}$ вариантов. Симметрично условие «Саша набрал $100$ до первой победы Оли» даёт столько же.
\par Дважды посчитаны последовательности, в которых и Паша, и Саша набрали $100$ до первой победы Оли. В них сначала идут $100$ букв П и $100$ букв С в произвольном порядке, а затем $100$ букв О. Таких последовательностей $\binom{200}{100}$.
\par Дописывая последнюю О, получаем в точности все требуемые игры: никто из участников не набирает $101$ очко раньше неё. По формуле включений-исключений ответ равен
\[
2\binom{300}{100}-\binom{200}{100}.
\]\par

\item Дано натуральное $n$. Докажите, что числа от $1$ до $2^n$ можно покрасить в два цвета так, чтобы не было одноцветной арифметической прогрессии длины $2n$. Разность прогрессии должна быть положительной.
\par\smallskip\noindent\textbf{Решение.} Положим $N=2^n$. Всего раскрасок множества $\{1,\ldots,N\}$ в два цвета $2^N$. Назовём раскраску плохой, если в ней есть одноцветная арифметическая прогрессия длины $2n$.
\par Такая прогрессия задаётся первым членом $a$ и положительной разностью $d$, причём
\[
a+(2n-1)d\le N.
\]
Поэтому число прогрессий равно
\[
M=\sum_{d=1}^{D}\bigl(N-(2n-1)d\bigr),
\qquad D=\left\lfloor\frac{N-1}{2n-1}\right\rfloor.
\]
Оценим его сверху:
\[
M\le\sum_{d=1}^{N-1}(N-d)
=\frac{N(N-1)}2
<\frac{N^2}{2}=2^{2n-1}.
\]
\par Для фиксированной прогрессии существует $2^{N-2n+1}$ раскрасок, в которых она одноцветна: выбираем один из двух цветов для её $2n$ членов, а остальные $N-2n$ чисел красим произвольно. Здесь $N=2^n\ge2n$ при $n\ge1$.
\par Считая плохие раскраски по всем прогрессиям, некоторые можем посчитать несколько раз, поэтому их число не больше
\[
M\,2^{N-2n+1}<2^{2n-1}\,2^{N-2n+1}=2^N.
\]
Плохих раскрасок строго меньше, чем всех. Значит, существует раскраска без одноцветной прогрессии указанной длины.\par

\item Даны $6000$ карточек, на которых написаны натуральные числа от $1$ до $6000$, каждое ровно один раз. Сколькими способами можно выбрать две карточки, чтобы сумма написанных на них чисел делилась на $100$? Порядок выбранных карточек не учитывается.
\par\smallskip\noindent\textbf{Решение.} В каждом классе остатков по модулю $100$ находится ровно $60$ карточек. Сначала посчитаем упорядоченные пары различных карточек.
\par Если остаток первой карточки не равен $0$ или $50$, подходящих вторых карточек ровно $60$: они принадлежат другому классу остатков. Первую карточку такого вида можно выбрать $6000-120=5880$ способами.
\par Если остаток первой карточки равен $0$ или $50$, вторая должна иметь тот же остаток. Исключая саму первую карточку, получаем $59$ вариантов. Первых карточек такого вида $120$.
\par Каждая неупорядоченная пара посчитана дважды, поэтому ответ равен
\[
\frac{5880\cdot60+120\cdot59}{2}=179\,940.
\]\par

\item Сколько семибуквенных слов можно составить из букв $a$ и $b$ так, чтобы непосредственно после каждой буквы $a$ стояла хотя бы одна буква $b$? Букве $a$ разрешается быть последней.
\par\smallskip\noindent\textbf{Решение.} Условие равносильно отсутствию двух букв $a$ подряд: любая не последняя буква $a$ должна сопровождаться буквой $b$. Переберём число $k$ букв $a$. Оно может принимать значения $0,1,2,3,4$.
\par Пусть позиции этих букв равны $1\le i_1<\cdots<i_k\le7$. Необходимые и достаточные ограничения --- $i_{j+1}\ge i_j+2$. Заменой
\[
j_t=i_t-(t-1)
\]
получаем выбор $k$ различных позиций из $8-k$. Поэтому для данного $k$ слов $\binom{8-k}{k}$. Суммируя, получаем
\[
\sum_{k=0}^{4}\binom{8-k}{k}=1+7+15+10+1=34.
\]\par

\item Петя хочет выписать все последовательности из $100$ натуральных чисел, в каждой из которых хотя бы раз встречается число $4$ или $5$, а любые два соседних члена отличаются не больше чем на $2$. Сколько последовательностей ему придётся выписать?
\par\smallskip\noindent\textbf{Решение.} Сначала рассмотрим последовательности длины $n$, минимум которых принадлежит множеству $\{1,2,3,4,5\}$. Шаблон последовательности задаётся её приращениями
\[
d_i=a_{i+1}-a_i\in\{-2,-1,0,1,2\},\qquad i=1,\ldots,n-1.
\]
Таких шаблонов $5^{n-1}$. Зафиксируем один из них и построим вспомогательную последовательность $b_1=0$, $b_{i+1}=b_i+d_i$. Если её минимум равен $m$, то для каждого $k\in\{1,2,3,4,5\}$ последовательность
\[
a_i=b_i-m+k
\]
имеет минимум $k$. Это единственная последовательность с данным шаблоном и данным минимумом. Следовательно, всего рассматриваемых последовательностей $5\cdot5^{n-1}=5^n$.
\par Любая последовательность из условия входит в это число, поскольку её минимум не больше $5$. Вычтем последовательности, в которых нет ни $4$, ни $5$. Их минимум не больше $3$. Если бы среди их членов встретилось число не меньше $6$, то между членом, не превосходящим $3$, и этим членом пришлось бы перейти через $4$ или $5$: перескочить сразу от числа не больше $3$ к числу не меньше $6$ при шаге не больше $2$ невозможно. Значит, все члены такой последовательности принадлежат $\{1,2,3\}$.
\par Обратно, любая последовательность из чисел $1,2,3$ имеет соседние разности не больше $2$ и не содержит ни $4$, ни $5$. Таких последовательностей $3^n$. Поэтому при $n=100$ ответ равен
\[
5^{100}-3^{100}.
\]\par

\item Сколькими способами можно представить $10^{10}$ в виде произведения трёх натуральных множителей, если произведения, отличающиеся порядком множителей, считаются различными?
\par\smallskip\noindent\textbf{Решение.} Поскольку $10^{10}=2^{10}5^{10}$, каждый из трёх множителей имеет вид
\[
x=2^{a_1}5^{b_1},\qquad y=2^{a_2}5^{b_2},\qquad z=2^{a_3}5^{b_3},
\]
где показатели --- неотрицательные целые числа. Равенство $xyz=10^{10}$ равносильно системе
\[
a_1+a_2+a_3=10,\qquad b_1+b_2+b_3=10.
\]
По формуле шаров и перегородок каждое уравнение имеет $\binom{12}{2}=66$ решений. Распределения показателей двойки и пятёрки независимы, а каждой их паре соответствует ровно одна упорядоченная тройка множителей. Поэтому ответ равен
\[
\binom{12}{2}^{\!2}=66^2=4356.
\]\par

\item Сколько существует $23$-значных чисел, сумма цифр которых равна восьми?
\par\smallskip\noindent\textbf{Решение.} Разобьём сумму $8$ на ненулевые цифры числа. Пусть цифра $d$ встречается $c_d$ раз, где $d=1,\ldots,8$. Тогда
\[
c_1+2c_2+\cdots+8c_8=8,
\qquad k=c_1+\cdots+c_8.
\]
Оставшиеся $23-k$ цифр --- нули. Для данного набора цифр всех перестановок
\[
\frac{23!}{(23-k)!\prod_{d=1}^{8}c_d!},
\]
но нужно исключить начинающиеся с нуля. Их
\[
\frac{22!}{(22-k)!\prod_{d=1}^{8}c_d!}.
\]
Разность равна
\[
\frac{k\,22!}{(23-k)!\prod_{d=1}^{8}c_d!}.
\]
Суммируем это выражение по всем разбиениям $8$ на ненулевые цифры.
\par Для вычисления сгруппируем варианты по числу $k$ ненулевых цифр, где $1\le k\le8$. Первая позиция обязательно ненулевая, а остальные $k-1$ ненулевых позиций выбираются $\binom{22}{k-1}$ способами. Упорядоченных наборов $k$ положительных цифр с суммой $8$ по формуле шаров и перегородок $\binom7{k-1}$. Верхняя граница цифры не мешает: каждая автоматически не больше $8$.
\par Итак, сумма равна
\[
\sum_{k=1}^{8}\binom{22}{k-1}\binom7{k-1}
=\sum_{j=0}^{7}\binom{22}{j}\binom7{7-j}
=\binom{29}{7}
=1\,560\,780,
\]
где использована формула Вандермонда.\par

\item Изначально все клетки доски $8\times8$ белые. За один ход можно поменять на противоположный цвет все клетки любого квадрата $3\times3$, все клетки любого квадрата $4\times4$ или одну из клеток $a1$ и $b2$. Любую ли чёрно-белую раскраску доски можно получить такими ходами? Обозначения клеток - обычные шахматные.
\par\smallskip\noindent\textbf{Решение.} Нет. Квадратов $3\times3$ на доске $(8-3+1)^2=36$, а квадратов $4\times4$ --- $(8-4+1)^2=25$. Кроме того, разрешены два отдельных хода с клетками $a1$ и $b2$. Итого имеется $36+25+2=63$ вида ходов.
\par Порядок ходов не влияет на итоговый цвет каждой клетки: он определяется чётностью числа её перекрашиваний. Два одинаковых хода взаимно уничтожаются, поэтому результат любой последовательности ходов можно получить, применив каждый вид хода не более одного раза. Выбрать набор применяемых ходов можно $2^{63}$ способами, следовательно, достижимых раскрасок не больше $2^{63}$.
\par Всего же чёрно-белых раскрасок доски $2^{64}$. Поскольку $2^{63}<2^{64}$, некоторые раскраски получить невозможно.\par

\item Школьники участвовали в олимпиаде, проходившей в два тура. В каждый из двух дней они рассаживались по нескольким кабинетам, причём никто не сидел в кабинете в одиночестве. Количество кабинетов в разные дни могло отличаться. Докажите, что найдутся два школьника $A$ и $B$ такие, что в первый день они писали олимпиаду в кабинетах с одинаковым количеством участников и во второй день также писали в кабинетах с одинаковым количеством участников. Сами кабинеты у этих школьников могли быть разными.
\par\smallskip\noindent\textbf{Решение.} Предположим, что требуемой пары школьников нет. Обозначим размеры кабинетов в первый день через $a_1\ge\cdots\ge a_k\ge2$, а во второй день --- через $b_1\ge\cdots\ge b_t\ge2$. Переставив дни, если необходимо, будем считать, что $a_1\ge b_1$.
\par Рассмотрим $a_1$ учеников самого большого кабинета первого дня. Во второй день они должны сидеть в кабинетах попарно различных размеров: иначе у двух из них совпадут размеры кабинетов в оба дня, вопреки предположению. Однако возможные размеры кабинетов второго дня --- только $2,3,\ldots,b_1$. Их всего $b_1-1\le a_1-1$, а учеников $a_1$. По принципу Дирихле два ученика окажутся в кабинетах одинакового размера. Полученное противоречие доказывает утверждение.\par

\item Докажите, что количество разбиений числа $n$ на натуральные слагаемые равно количеству разбиений числа $2n$ на ровно $n$ натуральных слагаемых.
\par\smallskip\noindent\textbf{Решение.} Пусть $n=\lambda_1+\cdots+\lambda_r$ --- разбиение на положительные слагаемые в невозрастающем порядке. Поскольку $r\le n$, дополним список нулями до длины $n$ и прибавим к каждому слагаемому единицу. Получится разбиение $2n$ на ровно $n$ натуральных слагаемых:
\[
(\lambda_1+1)+\cdots+(\lambda_r+1)
+\underbrace{1+\cdots+1}_{n-r}=2n.
\]
На диаграмме Юнга со слагаемыми-столбцами это добавление снизу строки из $n$ клеток.
\par Обратно, из каждого из $n$ слагаемых разбиения числа $2n$ вычтем $1$ и удалим нули. Получим разбиение $n$. Операции взаимно обратны, значит, задают биекцию.\par

\item Число называется \emph{шикарным}, если у него есть два делителя, отличающихся на $2$, и \emph{фееричным}, если у него есть два делителя, отличающихся на $3$. Каких чисел больше от $1$ до $6000$ включительно - шикарных или фееричных?
\par\smallskip\noindent\textbf{Решение.} Если два делителя числа отличаются на $3$, один из них чётный. Значит, само число чётное. Поэтому фееричных чисел от $1$ до $6000$ не больше $3000$.
\par Все числа, кратные $3$, шикарны благодаря делителям $1$ и $3$. Все числа, кратные $4$, также шикарны благодаря делителям $2$ и $4$. Чисел, делящихся хотя бы на одно из чисел $3$ и $4$, ровно
\[
\left\lfloor\frac{6000}{3}\right\rfloor
+\left\lfloor\frac{6000}{4}\right\rfloor
-\left\lfloor\frac{6000}{12}\right\rfloor
=2000+1500-500=3000.
\]
Кроме того, $35$ не кратно ни $3$, ни $4$, но шикарно: его делители $5$ и $7$ отличаются на $2$. Значит, шикарных чисел хотя бы $3001$, а фееричных не больше $3000$. Следовательно, \textbf{шикарных чисел больше}.\par

\item Докажите, что при любом натуральном $n$ уравнения
\[
x^2+y^2=n
\qquad\text{и}\qquad
x^2+y^2=2n
\]
имеют одинаковое количество решений в целых числах.
\par\smallskip\noindent\textbf{Решение.} Решению $(x,y)$ уравнения $x^2+y^2=n$ сопоставим пару
\[
(u,v)=(x+y,x-y).
\]
Тогда
\[
u^2+v^2=(x+y)^2+(x-y)^2=2(x^2+y^2)=2n.
\]
\par Обратно, пусть $u^2+v^2=2n$. Числа $u$ и $v$ одной чётности, поскольку квадраты имеют ту же чётность, что и сами числа, а сумма квадратов чётна. Поэтому полусумма и полуразность
\[
x=\frac{u+v}{2},\qquad y=\frac{u-v}{2}
\]
целые. При этом
\[
x^2+y^2=\frac{(u+v)^2+(u-v)^2}{4}
=\frac{u^2+v^2}{2}=n.
\]
Два преобразования взаимно обратны и задают биекцию между множествами целочисленных решений.\par

\item Рассмотрим всевозможные наборы из $100$ неотрицательных целых чисел, расположенных в неубывающем порядке и не превосходящих $100$, в которых сумма всех чисел кратна $10$. Докажите, что ровно половина этих наборов заканчивается на $100$.
\par\smallskip\noindent\textbf{Решение.} Изобразим набор $0\le a_1\le\cdots\le a_{100}\le100$ диаграммой Юнга: столбцы имеют высоты $a_1,\ldots,a_{100}$ и выровнены по нижней стороне квадрата $100\times100$. Столбцы нулевой высоты также учитываем. Возьмём дополнение диаграммы до всего квадрата и прочитаем его сбоку: новыми слагаемыми будут длины строк дополнения, снизу вверх.
\par Формально новый набор задаётся числами
\[
b_j=\#\{i:a_i<j\},\qquad j=1,\ldots,100.
\]
Это неубывающий набор из $100$ целых чисел от $0$ до $100$. Сумма новых слагаемых равна площади дополнения:
\[
\sum_{j=1}^{100}b_j=10000-\sum_{i=1}^{100}a_i,
\]
поэтому она также кратна $10$.
\par Преобразование взаимно обратно самому себе. Действительно, поскольку набор $a$ неубывающий, неравенство $b_j<i$ равносильно $a_i\ge j$. Следовательно, при повторном преобразовании получаем
\[
c_i=\#\{j:b_j<i\}=\#\{j:1\le j\le a_i\}=a_i.
\]
\par Осталось проверить, что два нужных класса меняются местами. Если $a_{100}=100$, в исходной диаграмме есть полный столбец. Тогда каждая строка дополнения содержит не больше $99$ клеток, поэтому $b_{100}<100$. Если же $a_{100}<100$, верхняя строка исходной диаграммы пуста, поэтому верхняя строка дополнения полная и $b_{100}=100$.
\par Таким образом, построена биекция между допустимыми наборами, заканчивающимися на $100$, и всеми остальными допустимыми наборами. Их количества равны, значит, ровно половина наборов заканчивается на $100$.\par

\item Докажите, что количество разбиений числа $n$ в сумму не более чем $k$ слагаемых, каждое из которых не превосходит $\ell$, равно количеству разбиений числа $n$ в сумму не более чем $\ell$ слагаемых, каждое из которых не превосходит $k$.
\par\smallskip\noindent\textbf{Решение.} Изобразим разбиение диаграммой Юнга, в которой каждому слагаемому соответствует столбец. Ограничения означают, что столбцов не больше $k$, а высота каждого не больше $\ell$.
\par Транспонируем диаграмму. Теперь столбцов не больше $\ell$, а их высоты не больше $k$. Число клеток, то есть сумма слагаемых $n$, сохраняется. Повторное транспонирование возвращает исходную диаграмму, следовательно, это требуемая биекция.\par

\item Сколько существует способов разбить число $2015$ на приблизительно равные натуральные слагаемые? Слагаемых может быть одно или несколько. Числа называются приблизительно равными, если их разность не больше $1$. Например:
\[
\underbrace{1+1+\ldots+1}_{2015\text{ слагаемых}},
\qquad 671+672+672,
\qquad 1007+1008.
\]
\par\smallskip\noindent\textbf{Решение.} Зафиксируем число слагаемых $\ell$, где $1\le\ell\le2015$. Разделим $2015$ на $\ell$ с остатком:
\[
2015=q\ell+r,\qquad 0\le r<\ell.
\]
Приблизительно равные слагаемые могут принимать только два соседних значения. Поэтому единственное разбиение с $\ell$ слагаемыми состоит из $r$ слагаемых $q+1$ и $\ell-r$ слагаемых $q$. Все они натуральные, поскольку $q\ge1$.
\par Таким образом, каждому $\ell=1,\ldots,2015$ соответствует ровно одно разбиение. Разные $\ell$ дают разные разбиения, поэтому ответ:
\[
2015.
\]\par

\item Назовём билет с номером от $000000$ до $999999$ отличным, если некоторые две соседние цифры его номера отличаются на $5$. Найдите число отличных билетов. Нули в начале номера допускаются.
\par\smallskip\noindent\textbf{Решение.} Зафиксируем первую цифру номера. Всего билетов с такой первой цифрой $10^5$. Посчитаем среди них неотличные, то есть номера без соседних цифр, отличающихся на $5$.
\par У любой цифры $a$ есть ровно одна цифра, отличающаяся от неё на $5$: это $a+5$ при $a\le4$ и $a-5$ при $a\ge5$. Поэтому для каждой следующей позиции запрещена ровно одна цифра и разрешены остальные девять. При фиксированной первой цифре неотличных билетов $9^5$, а отличных --- $10^5-9^5$.
\par Первую цифру можно выбрать десятью способами. Ноль ничем не отличается от остальных вариантов, поскольку ведущие нули разрешены. Следовательно, ответ равен
\[
10(10^5-9^5)=409\,510.
\]\par

\end{problems}
\end{document}
